\tinysidebar{\debug{exercises/{test/question.tex}}}
\begin{ex}
Get an integer n from the user and then compute the sum of the leftmost 4 digits of n. Do this with repetition code.

\protect\hypertarget{anchor-1}Solution.

\begin{consolethree}
#include <iostream>

int main()
{
    int n;
    std::cin >> n;
    int s = 0;
    s = s + n % 10;
    n = n / 10;
    s = s + n % 10;
    n = n / 10;
    s = s + n % 10;
    n = n / 10;
    s = s + n % 10;
    n = n / 10;
    std::cout << s << ", " << n;
    return 0;
} 


\begin{ex} Get integer \texttt{x}, \texttt{y} from the user and print \texttt{x}, \texttt{y}, except that if \texttt{x }is negative you print \texttt{0} instead of \texttt{x} and if \texttt{y} is negative you print \texttt{0} instead of \texttt{y}. Here's a test case

\begin{console}
2 3

2 3 
\end{console}

Here's another test case

\begin{console}
-1 3

0 3 
\end{console}

and another test:

\begin{console}
2 -2

2 0 
\end{console}

and finally

\begin{console}
-5 -2

0 0 
\end{console

You can write your program with boolean conditions where each boolean condition involves both \texttt{x} and \texttt{y}. 
\end{ex}
